% Part: normal-modal-logic
% Chapter: tableaux
% Section: countermodels

\documentclass[../../../include/open-logic-section]{subfiles}

\begin{document}

\olfileid[hi]{nml}{tab}{cou}
      
\olsection{\usetoken{P}{tableau} से प्रतिप्रतिरूप}

\begin{explain}
  पूर्णता प्रमेय का प्रमाण केवल यह नहीं दिखाता कि यदि $\Entails !A$, तो
  $\Proves !A$; यदि $\Entails/ A$, तो वह हमें $!A$ के प्रतिप्रतिरूप बनाने
  की विधि भी देता है। $\Log{K}$ की स्थिति में यह विधि एक \emph{निर्णय
    प्रक्रिया} बनाती है। मान लें $\Entails/ !A$। तब
  \olref[cpl]{prop:complete-tableau} का प्रमाण पूर्ण !!{tableau} बनाने की
  विधि देता है। वास्तव में यह विधि सदा समाप्त होती है। $\Log{K}$ के
  प्रतिज्ञप्तिक नियम केवल कम जटिलता वाले उपसर्गयुक्त !!{formula} जोड़ते हैं,
  अर्थात् किसी भी चिह्नित सूत्र $\sFmla{S}{!A}[\sigma]$ के लिए प्रत्येक
  प्रतिज्ञप्तिक नियम किसी शाखा पर केवल एक बार लगाना आवश्यक है। नए उपसर्ग
  केवल
  \iftag{prvBox}{$\TRule{\False}{\Box}$}{}\iftag{notprvBox,notprvDiamond}{}{
    और }\iftag{prvDiamond}{$\TRule{\True}{\Diamond}$}{}
  \iftag{notprvBox,notprvDiamond}{नियम}{नियमों} से उत्पन्न होते हैं; और
  \iftag{notprvBox,notprvDiamond}{उसे}{उन्हें} भी केवल एक बार लगाना होता
  है (और एक ही नया उपसर्ग उत्पन्न होता है)।
  \iftag{prvBox}{$\TRule{\True}{\Box}$}{}\iftag{notprvBox,notprvDiamond}{}{
    और }\iftag{prvDiamond}{$\TRule{\False}{\Diamond}$}{} को संभवतः अनेक
  बार लगाना होता है, किंतु प्रत्येक उपसर्ग के लिए केवल एक बार; और केवल
  परिमिततः अनेक नए उपसर्ग उत्पन्न होते हैं। अतः निर्माण परिमिततः अनेक
  चरणों के बाद या तो संवृत शाखा देता है, या पूर्ण शाखा।

  खुली पूर्ण शाखा वाला टैब्लो बन जाने पर
  \olref[cpl]{thm:tableau-completeness} का प्रमाण हमें ऐसा स्पष्ट प्रतिरूप
  देता है जो उपसर्गयुक्त !!{formula}ों के मूल समुच्चय को संतुष्ट करता है।
  अतः केवल यही नहीं कि यदि $\Gamma \Entails !A$, तो कोई संवृत !!{tableau}
  है और $\Gamma \Proves !A$; यदि हम संवृत !!{tableau} उचित ढंग से खोजें और
  अंत में कोई ``पूर्ण'' !!{tableau} पाएँ, तो हमें केवल $\Gamma \Entails/
  !A$ का ज्ञान ही नहीं होगा, बल्कि हम वास्तव में प्रतिप्रतिरूप बना सकेंगे।
\end{explain}

\iftag{prvBox}{
\begin{ex}
  हम जानते हैं कि $\Proves/ \Box(p \lor q) \lif (\Box p \lor \Box
  q)$। टैब्लो का निर्माण इस प्रकार आरंभ होता है:
  \begin{oltableau}
    [\pFmla{\False}{\Box(p \lor q) \lif (\Box p \lor \Box q)}{1},
      just = \TAss, checked
      [\pFmla{\True}{\Box(p \lor q)}{1},
        just = {\TRule{\False}{\lif}[1]}, 
        [\pFmla{\False}{\Box p \lor \Box q}{1},
          just = {\TRule{\False}{\lif}[1]}, checked
          [\pFmla{\False}{\Box p}{1},
            just = {\TRule{\False}{\lor}[3]}, checked
            [\pFmla{\False}{\Box q}{1},
              just = {\TRule{\False}{\lor}[3]}, checked
              [\pFmla{\False}{p}{1.1}, 
                just = {\TRule{\False}{\Box}[4]}, checked
                [\pFmla{\False}{q}{1.2}, 
                  just = {\TRule{\False}{\Box}[5]}, checked
                ]
              ]
            ]
          ]
        ]
      ]
    ]
  \end{oltableau}
  निस्संदेह !!{tableau} अभी पूरा नहीं हुआ है। अगले चरण में हम बिना सही-चिह्न
  वाली एकमात्र पंक्ति पर विचार करते हैं: पंक्ति~$2$ का उपसर्गयुक्त
  !!{formula} $\sFmla{\True}{\Box(p \lor q)}[1]$। संवृत टैब्लो का निर्माण
  कहता है कि शाखा पर प्रयुक्त प्रत्येक उपसर्ग, अर्थात् $1.1$ और $1.2$
  दोनों, के लिए $\TRule{\True}{\Box}$ नियम लगाएँ:
  \begin{oltableau}
    [\pFmla{\False}{\Box(p \lor q) \lif (\Box p \lor \Box q)}{1},
      just = \TAss, checked
      [\pFmla{\True}{\Box(p \lor q)}{1},
        just = {\TRule{\False}{\lif}[1]}, 
        [\pFmla{\False}{\Box p \lor \Box q}{1},
          just = {\TRule{\False}{\lif}[1]}, checked
          [\pFmla{\False}{\Box p}{1},
            just = {\TRule{\False}{\lor}[3]}, checked
            [\pFmla{\False}{\Box q}{1},
              just = {\TRule{\False}{\lor}[3]}, checked
              [\pFmla{\False}{p}{1.1}, 
                just = {\TRule{\False}{\Box}[4]}, checked
                [\pFmla{\False}{q}{1.2}, 
                  just = {\TRule{\False}{\Box}[5]}, checked
                  [\pFmla{\True}{p \lor q}{1.1}, 
                    just = {\TRule{\True}{\Box}[2]}
                    [\pFmla{\True}{p \lor q}{1.2}, 
                      just = {\TRule{\True}{\Box}[2]}
                    ]
                  ]
                ]
              ]
            ]
          ]
        ]
      ]
    ]
  \end{oltableau}
  अब पंक्तियों 2, 8 और 9 पर सही-चिह्न नहीं हैं। किंतु कोई नया उपसर्ग नहीं
  जोड़ा गया है, इसलिए हम सभी परिणामी शाखाओं पर (जब तक वे संवृत न हों)
  पंक्तियों~8 और~9 को $\TRule{\True}{\lor}$ लगाते हैं:
  \begin{oltableau}
    [\pFmla{\False}{\Box(p \lor q) \lif (\Box p \lor \Box q)}{1},
      just = \TAss, checked
      [\pFmla{\True}{\Box(p \lor q)}{1},
        just = {\TRule{\False}{\lif}[1]}, checked
        [\pFmla{\False}{\Box p \lor \Box q}{1},
          just = {\TRule{\False}{\lif}[1]}, checked
          [\pFmla{\False}{\Box p}{1},
            just = {\TRule{\False}{\lor}[3]}, checked
            [\pFmla{\False}{\Box q}{1},
              just = {\TRule{\False}{\lor}[3]}, checked
              [\pFmla{\False}{p}{1.1}, 
                just = {\TRule{\False}{\Box}[4]}, checked
                [\pFmla{\False}{q}{1.2}, 
                  just = {\TRule{\False}{\Box}[5]}, checked
                  [\pFmla{\True}{p \lor q}{1.1}, 
                    just = {\TRule{\True}{\Box}[2]}, checked
                    [\pFmla{\True}{p \lor q}{1.2}, 
                      just = {\TRule{\True}{\Box}[2]}, checked
                      [\pFmla{\True}{p}{1.1},
                        just = {\TRule{\True}{\lor}[8]}, checked, close
                      ]
                      [\pFmla{\True}{q}{1.1},
                        just = {\TRule{\True}{\lor}[8]}, checked
                        [\pFmla{\True}{p}{1.2},
                          just = {\TRule{\True}{\lor}[9]}, checked]
                        [\pFmla{\True}{q}{1.2},
                          just = {\TRule{\True}{\lor}[9]}, checked, close]
                      ]
                    ]
                  ]
                ]
              ]
            ]
          ]
        ]
      ]
    ]
  \end{oltableau}
  एक खुली शाखा शेष है और वह पूर्ण है। उससे हम ऐसा प्रतिरूप परिभाषित करते
  हैं जिसके जगत $W = \{1, 1.1, 1.2\}$ (खुली शाखा पर आने वाले एकमात्र
  उपसर्ग), अभिगम्यता संबंध $R = \{\tuple{1, 1.1}, \tuple{1, 1.2}\}$,
  तथा नियतन $V(p) = \{1.2\}$ (क्योंकि पंक्ति~11 में
  $\sFmla{\True}{p}[1.2]$ है) और $V(q) = \{1.1\}$ (क्योंकि पंक्ति~10
  में $\sFmla{\True}{q}[1.1]$ है) हैं। प्रतिरूप को
  \olref{fig:counter-Box} में चित्रित किया गया है, और आप सत्यापित कर सकते
  हैं कि वह $\Box(p \lor q) \lif (\Box p \lor \Box q)$ का प्रतिप्रतिरूप
  है।
  \begin{figure}
  \begin{center}
    \begin{tikzpicture}[modal]
      \node[world] (w1) [label={[align=right]right:\mFalse{p}\\ \mFalse{q}}]{$1$}; 
      \node[world] (w2) [label={[align=right]right:\mFalse{p}\\ \mTrue{q}},
        above left=of w1]{$1.1$}; 
      \node[world] (w3) [label={[align=right]right:\mTrue{p}\\ \mFalse{q}},
        above right=of w1] {$1.2$};
      \draw[->] (w1) to (w2);
      \draw[->] (w1) to (w3);
    \end{tikzpicture}
  \end{center}
  \caption{$\Box(p \lor q) \lif (\Box p \lor \Box q)$ का प्रतिप्रतिरूप।}
  \ollabel{fig:counter-Box}
\end{figure}
\end{ex}          
}
{
\begin{ex}
  हम जानते हैं कि $\Proves/ (\Diamond p \land \Diamond q) \lif
  \Diamond(p \land q)$। टैब्लो का निर्माण इस प्रकार आरंभ होता है:
  \begin{oltableau}
    [\pFmla{\False}{(\Diamond p \land \Diamond q) \lif \Diamond(p \land q)}{1},
      just = \TAss, checked
      [\pFmla{\True}{\Diamond p \land \Diamond q}{1},
        just = {\TRule{\False}{\lif}[1]}, checked
        [\pFmla{\False}{\Diamond(p \land q)}{1},
          just = {\TRule{\False}{\lif}[1]}
          [\pFmla{\True}{\Diamond p}{1},
            just = {\TRule{\True}{\land}[2]}, checked
            [\pFmla{\True}{\Diamond q}{1},
              just = {\TRule{\True}{\land}[2]}, checked
              [\pFmla{\True}{p}{1.1}, 
                just = {\TRule{\True}{\Diamond}[4]}, checked
                [\pFmla{\True}{q}{1.2}, 
                  just = {\TRule{\True}{\Diamond}[5]}, checked
                ]
              ]
            ]
          ]
        ]
      ]
    ]
  \end{oltableau}
  निस्संदेह !!{tableau} अभी पूरा नहीं हुआ है। अगले चरण में हम बिना सही-चिह्न
  वाली एकमात्र पंक्ति पर विचार करते हैं: पंक्ति~$3$ का उपसर्गयुक्त
  !!{formula} $\sFmla{\True}{\Diamond(p \land q)}[1]$। संवृत टैब्लो का
  निर्माण कहता है कि शाखा पर प्रयुक्त प्रत्येक उपसर्ग, अर्थात् $1.1$ और
  $1.2$ दोनों, के लिए $\TRule{\True}{\Diamond}$ नियम लगाएँ:
  \begin{oltableau}
    [\pFmla{\False}{\Diamond(p \land q) \lif (\Diamond p \land \Diamond q)}{1},
      just = \TAss, checked
      [\pFmla{\True}{\Diamond p \land \Diamond q}{1},
        just = {\TRule{\False}{\lif}[1]}, checked
        [\pFmla{\False}{\Diamond (p \land q)}{1},
          just = {\TRule{\False}{\lif}[1]}
          [\pFmla{\True}{\Diamond p}{1},
            just = {\TRule{\True}{\land}[2]}, checked
            [\pFmla{\True}{\Diamond q}{1},
              just = {\TRule{\True}{\land}[2]}, checked
              [\pFmla{\True}{p}{1.1}, 
                just = {\TRule{\True}{\Diamond}[4]}, checked
                [\pFmla{\True}{q}{1.2}, 
                  just = {\TRule{\True}{\Diamond}[5]}, checked
                  [\pFmla{\False}{p \land q}{1.1}, 
                    just = {\TRule{\False}{\Diamond}[3]}
                    [\pFmla{\False}{p \land q}{1.2}, 
                      just = {\TRule{\False}{\Diamond}[3]}
                    ]
                  ]
                ]
              ]
            ]
          ]
        ]
      ]
    ]
  \end{oltableau}
  अब पंक्तियों 3, 8 और 9 पर सही-चिह्न नहीं हैं। किंतु कोई नया उपसर्ग नहीं
  जोड़ा गया है, इसलिए हम सभी परिणामी शाखाओं पर (जब तक वे संवृत न हों)
  पंक्तियों~8 और~9 को $\TRule{\False}{\land}$ लगाते हैं:
  \begin{oltableau}
    [\pFmla{\False}{(\Diamond p \land \Diamond q) \lif \Diamond(p \land q)}{1},
      just = \TAss, checked
      [\pFmla{\True}{\Diamond p \land \Diamond q}{1},
        just = {\TRule{\False}{\lif}[1]}, checked
        [\pFmla{\False}{\Diamond(p \land q)}{1},
          just = {\TRule{\False}{\lif}[1]}
          [\pFmla{\True}{\Diamond p}{1},
            just = {\TRule{\True}{\land}[2]}, checked
            [\pFmla{\True}{\Diamond q}{1},
              just = {\TRule{\True}{\land}[2]}, checked
              [\pFmla{\True}{p}{1.1}, 
                just = {\TRule{\True}{\Diamond}[4]}, checked
                [\pFmla{\True}{q}{1.2}, 
                  just = {\TRule{\True}{\Diamond}[5]}, checked
                  [\pFmla{\False}{p \land q}{1.1}, 
                    just = {\TRule{\False}{\Diamond}[3]}, checked
                    [\pFmla{\False}{p \land q}{1.2}, 
                      just = {\TRule{\False}{\Diamond}[3]}, checked
                      [\pFmla{\False}{p}{1.1},
                        just = {\TRule{\False}{\land}[8]}, checked, close
                      ]
                      [\pFmla{\False}{q}{1.1},
                        just = {\TRule{\False}{\land}[8]}, checked
                        [\pFmla{\False}{p}{1.2},
                          just = {\TRule{\False}{\land}[9]}, checked]
                        [\pFmla{\False}{q}{1.2},
                          just = {\TRule{\False}{\land}[9]}, checked, close]
                      ]
                    ]
                  ]
                ]
              ]
            ]
          ]
        ]
      ]
    ]
  \end{oltableau}
  एक खुली शाखा शेष है और वह पूर्ण है। उससे हम ऐसा प्रतिरूप परिभाषित करते
  हैं जिसके जगत $W = \{1, 1.1, 1.2\}$ (खुली शाखा पर आने वाले एकमात्र
  उपसर्ग), अभिगम्यता संबंध $R = \{\tuple{1, 1.1}, \tuple{1, 1.2}\}$,
  तथा नियतन $V(p) = \{1.1\}$ (क्योंकि पंक्ति~6 में
  $\sFmla{\True}{p}[1.1]$ है) और $V(q) = \{1.2\}$ (क्योंकि पंक्ति~7
  में $\sFmla{\True}{q}[1.1]$ है) हैं। प्रतिरूप को
  \olref{fig:counter-Diamond} में चित्रित किया गया है, और आप सत्यापित कर
  सकते हैं कि वह $(\Diamond p \land \Diamond q) \lif \Diamond (p \land
  q)$ का प्रतिप्रतिरूप है।
  \begin{figure}
  \begin{center}
    \begin{tikzpicture}[modal]
      \node[world] (w1) [label={[align=right]right:\mFalse{p}\\ \mFalse{q}}]{$1$}; 
      \node[world] (w2) [label={[align=right]right:\mTrue{p}\\ \mFalse{q}},
        above left=of w1]{$1.1$}; 
      \node[world] (w3) [label={[align=right]right:\mFalse{p}\\ \mTrue{q}},
        above right=of w1] {$1.2$};
      \draw[->] (w1) to (w2);
      \draw[->] (w1) to (w3);
    \end{tikzpicture}
  \end{center}
  \caption{$(\Diamond p \land \Diamond q) \lif
    \Diamond (p \land q)$ का प्रतिप्रतिरूप।}
  \ollabel{fig:counter-Diamond}
\end{figure}
\end{ex}          
}

\end{document}
